% Imporditud: tools/import_eksperimendid.py
% Allikas: Eesti füüsikaolümpiaadi eksperimendi ülesannete
%   kogu (Rael Kalda, 2023), ülesanne Ü4, lahendus L4.
\setAuthor{EFO žürii}
\setRound{piirkonnavoor}
\setYear{1998}
\setNumber{P E2}
\setDifficulty{1}
\setTopic{dünaamika}
\setVahendid{kumminiit, statiiv, mõõtejoonlaud, tuntud massiga keha, tundmatu massiga keha}
\setPunktid{10}
\setAllikas{Eksperimendi kogu Ü4, lahendus lk 35}
\setLahendusseis{toores}

\prob{Keha mass}
Määra keha mass. Kirjeldada ülesande lahendust.

\hint

\solu
Teooria: Keha massi saab avaldada kehale mõjuvast raskusjõust. Raskusjõu valem on F = mg, kus F on raskusjõud, m on keha mass, g on konstant, mis arvuliselt on võrdne 10 N/kg (1 p). Kehale mõjuvat raskusjõudu saab mõõta dünamomeetriga. Dünamomeetriga mõõtes on tekkinud elastsusjõud suuruselt võrdne kehale mõjuva raskusjõuga (1 p). Dünamomeeter valmistatakse kumminiidist. Kumminiidi pikenemisel tekkiv elastsusjõud on võrdeline kumminiidi pikenemisega (1 p). Tuntud massiga kehale mõjuva raskusjõu F mõjul pikeneb kumminiit pikkuse s võrra. Tundmatu massiga kehale mõjuva raskusjõu $F_1$ mõjul pikeneb kumminiit pikkuse $s_1$ võrra. Tundmatu ja tuntud massiga kehale mõjuvate raskusjõudude jagatis on võrdne kumminiidi pikenemiste jagatisega $F_1$/F = $s_1$/s (1 p).

Mõõtmised: Dünamomeetri valmistamine. Teeme kumminiidi ühte otsa aasa ja kinnitame kumminiidi teise otsa statiivi klambri vahele (1 p). Mõõdame koormiseta kumminiidi pikkuse. Mõõdame kumminiidi pikkuse, kui selle otsa on riputatud tuntud massiga koormis. Mõõdame kumminiidi pikkuse, kui selle otsa on riputatud tundmatu massiga koormis (1 p).

Arvutused: Arvutame tuntud massiga koormisele mõjuva raskusjõu valemist F = mg. Arvutame kumminiidi pikenemise s tuntud massiga koormise mõjul (1 p). Arvutame kumminiidi pikenemise $s_1$ tundmatu massiga koormise mõjul (1 p). Arvutame tundmatu massiga kehale mõjuv raskusjõud $F_1$ = F $\cdot$ $s_1$/s (1 p). Arvutame tundmatu massiga keha massi $m_1$ = $F_1$/g (1 p).
\probend
